% Options for packages loaded elsewhere
\PassOptionsToPackage{unicode}{hyperref}
\PassOptionsToPackage{hyphens}{url}
\PassOptionsToPackage{dvipsnames,svgnames,x11names}{xcolor}
%
\documentclass[
  10pt,
  a4paper,
]{article}
\usepackage{amsmath,amssymb}
\usepackage{setspace}
\usepackage{iftex}
\ifPDFTeX
  \usepackage[T1]{fontenc}
  \usepackage[utf8]{inputenc}
  \usepackage{textcomp} % provide euro and other symbols
\else % if luatex or xetex
  \usepackage{unicode-math} % this also loads fontspec
  \defaultfontfeatures{Scale=MatchLowercase}
  \defaultfontfeatures[\rmfamily]{Ligatures=TeX,Scale=1}
\fi
\usepackage{lmodern}
\ifPDFTeX\else
  % xetex/luatex font selection
  \setmainfont[]{Liberation Serif}
  \setsansfont[]{Liberation Sans}
  \setmonofont[]{DejaVu Sans Mono}
\fi
% Use upquote if available, for straight quotes in verbatim environments
\IfFileExists{upquote.sty}{\usepackage{upquote}}{}
\IfFileExists{microtype.sty}{% use microtype if available
  \usepackage[]{microtype}
  \UseMicrotypeSet[protrusion]{basicmath} % disable protrusion for tt fonts
}{}
\makeatletter
\@ifundefined{KOMAClassName}{% if non-KOMA class
  \IfFileExists{parskip.sty}{%
    \usepackage{parskip}
  }{% else
    \setlength{\parindent}{0pt}
    \setlength{\parskip}{6pt plus 2pt minus 1pt}}
}{% if KOMA class
  \KOMAoptions{parskip=half}}
\makeatother
\usepackage{xcolor}
\usepackage[margin=22mm]{geometry}
\setlength{\emergencystretch}{3em} % prevent overfull lines
\providecommand{\tightlist}{%
  \setlength{\itemsep}{0pt}\setlength{\parskip}{0pt}}
\setcounter{secnumdepth}{-\maxdimen} % remove section numbering
\ifLuaTeX
\usepackage[bidi=basic]{babel}
\else
\usepackage[bidi=default]{babel}
\fi
\babelprovide[main,import]{korean}
\ifPDFTeX
\else
\babelfont{rm}[]{Liberation Serif}
\fi
% get rid of language-specific shorthands (see #6817):
\let\LanguageShortHands\languageshorthands
\def\languageshorthands#1{}
\usepackage{xeCJK}
\setCJKmainfont{Noto Serif CJK KR}
\setCJKsansfont{Noto Sans CJK KR}
\setCJKmonofont{Noto Sans Mono CJK KR}
\xeCJKsetup{CJKspace=true}
\usepackage{xurl}
\usepackage{fvextra}
\usepackage{fancyhdr}
\AtBeginDocument{\DefineVerbatimEnvironment{verbatim}{Verbatim}{breaklines=true,breakanywhere=true,fontsize=\small}}
\setlength{\emergencystretch}{4em}
\setlength{\headheight}{15pt}
\pagestyle{fancy}\fancyhf{}
\fancyhead[L]{\small 중성자별 구조와 비방사 진동}
\fancyfoot[C]{\thepage}
\renewcommand{\headrulewidth}{0.3pt}
\setcounter{tocdepth}{2}
\ifLuaTeX
  \usepackage{selnolig}  % disable illegal ligatures
\fi
\usepackage{bookmark}
\IfFileExists{xurl.sty}{\usepackage{xurl}}{} % add URL line breaks if available
\urlstyle{same}
\hypersetup{
  pdftitle={A19 · 통제 비교와 분해능 판단},
  pdflang={ko-KR},
  colorlinks=true,
  linkcolor={Maroon},
  filecolor={Maroon},
  citecolor={Blue},
  urlcolor={Blue},
  pdfcreator={LaTeX via pandoc}}

\title{A19 · 통제 비교와 분해능 판단}
\usepackage{etoolbox}
\makeatletter
\providecommand{\subtitle}[1]{% add subtitle to \maketitle
  \apptocmd{\@title}{\par {\large #1 \par}}{}{}
}
\makeatother
\subtitle{과제 큐시트}
\author{}
\date{}

\begin{document}
\maketitle

\setstretch{1.12}
\textbf{선행:} A07, A10, A18, A20　 \textbf{독립 작업:} 8--14시간
(계획치)

\subsection{차이가 일정하면 두 계산도 수렴한
것일까?}\label{uxcc28uxc774uxac00-uxc77cuxc815uxd558uxba74-uxb450-uxacc4uxc0b0uxb3c4-uxc218uxb834uxd55c-uxac83uxc77cuxae4c}

두 별의 주파수가 refinement마다 함께 이동하여 차이만 거의 일정하다. 이
경우 확인한 것은 차이인가, 각각의 주파수인가? 먼저 답을 분리하자.

먼저 한두 문장으로 답한 뒤 아래 설명을 읽자. 첫 예측이 틀렸거나 아직
모르겠다고 적는 것은 감점하지 않는다.

\subsection{필요한 관계와 읽을
위치}\label{uxd544uxc694uxd55c-uxad00uxacc4uxc640-uxc77duxc744-uxc704uxce58}

\(\delta_{num}=\max(|\Delta_1-\Delta_0|,|\Delta_2-\Delta_1|)\). 독립
root·roundoff 척도보다 큰 마지막 변화량이 증가하면
convergence\_unresolved이다.

\begin{itemize}
\tightlist
\item
  \path|notes/chapters/ch10.md|
\item
  \path|learner/src/diagnostics/ns_structure_diagnostics.f90|
\item
  \path|learner/src/diagnostics/ns_mode_diagnostics.f90|
\end{itemize}

\subsection{계산으로
옮기기}\label{uxacc4uxc0b0uxc73cuxb85c-uxc62euxae30uxae30}

분류를 먼저 시도한 뒤 제공 판정과 대조한다. 한 pair의 성공은 A21의 전체
수렴 검사를 대신하지 않는다.

\begin{enumerate}
\def\labelenumi{\arabic{enumi}.}
\tightlist
\item
  제공 구조·음향 진단과 차이 분류기를 호출한다. 직접 구현은 선택한다.
\item
  본문의 세 refinement 기록을 먼저 스스로 분류한 뒤 판정과 대조한다.
\item
  한 pair의 결과와 전체 A21 수렴 자료를 구분하여 해석한다.
\end{enumerate}

\newpage

\subsection{실행해서 무엇을 볼
것인가}\label{uxc2e4uxd589uxd574uxc11c-uxbb34uxc5c7uxc744-uxbcfc-uxac83uxc778uxac00}

패키지 루트에서 실행한다. \texttt{-\/-plan}은 경로만 확인하며 계산
성공을 뜻하지 않는다.

\begin{verbatim}
python tools/run_assignment.py A19 --plan
python tools/run_assignment.py A19 --run \
  --workspace learner --output ../A19_result.json
\end{verbatim}

\subsection{결과를
설명하기}\label{uxacb0uxacfcuxb97c-uxc124uxba85uxd558uxae30}

서로 다른 오류가 상쇄될 수 있다. 절대 주파수 수열과 paired difference를
함께 읽고, 아직 분해하지 못했다는 결론을 물리적 무효과로 바꾸지 않는다.

\textbf{남길 증거.} 핵심 자동시험; 세 수준의 차이·floor·판정 표; 같은
core/M/closure를 유지하는 계획서 또는 실제 pair 증거; mass residual의
주파수 기여 한계

\textbf{판단 기준.} 2점 예: 자동시험·기록 해석·실제 pair 생성의 증거
수준을 분리하고 unresolved 판정을 정확히 내린다. 최초 예측의 정오보다,
결과를 본 뒤 이유를 고치고 다른 입력에도 적용하는지를 본다.

\subsection{막혔을 때}\label{uxb9c9uxd614uxc744-uxb54c}

\textbf{힌트 1.} 각 refinement에서 baseline과 variant를 같은 설정으로
계산해야 한다.

\textbf{힌트 2.} 주어진 p1 변화량의 마지막 증가를 먼저 확인하라. fine
값이 크다는 이유로 resolved라 부르지 않는다.

\subsection{다음 만남에 다시
묻기}\label{uxb2e4uxc74c-uxb9ccuxb0a8uxc5d0-uxb2e4uxc2dc-uxbb3buxae30}

다음 회차에는 분해능 이하·수렴 미확인·모드 식별 미확인의 세 상황에 다음
실험을 각각 하나씩 고른다. 모든 경우에 tolerance만 줄인다고 답하지
않는다. 먼저 짧게 답한 뒤 본문과 자신의 기록을 열어 확인한다.

\textbf{범위.} 실패한 저장 배경 수열과 native EOS 함수를 직접 적분한
결과를 구별한다. 한 pair의 성공을 전체 오차 확인으로 읽지 않는다.

\end{document}
